% subject: physics
% type: assertion_reason
% has_diagram: false

\item \textbf{Assertion (A):} The work done by friction on a body sliding down a rough inclined plane is negative.

\textbf{Reason (R):} Friction always opposes the relative motion between two surfaces in contact.

\begin{tasks}(1)
    \task Both A and R are true and R is the correct explanation of A \ans
    \task Both A and R are true but R is not the correct explanation of A
    \task A is true but R is false
    \task A is false but R is true
\end{tasks}

\begin{solution}
\begin{align*}
\intertext{The assertion is true. When a body slides down a rough incline, friction acts up the incline (opposing motion). Since the displacement is down the incline, the angle between friction and displacement is $180^\circ$.}
W_f &= f \cdot d \cdot \cos 180^\circ \\
    &= -fd
\intertext{The work done is negative. The reason is also true — friction always opposes relative motion. This is precisely why friction acts opposite to the displacement, making the work negative. Hence R correctly explains A.}
\end{align*}
\end{solution}
